\section{A certified classical spiral-inflation fragment}
\label{sec:spiralinflation}

The triangle Finner test does not exhaust the restrictions imposed by source
independence. A small classical inflation supplies a stronger, still tractable
compatibility test. We use a fixed six-variable spiral construction from the
inflation framework of Wolfe, Spekkens, and Fritz
\cite{wolfe2019classicalinflation}; it is not a complete inflation hierarchy.
Its resource contract is explicitly classical fanout.

\subsection{Copied sources and observable marginal constraints}
Let $U,V,W$ be independent classical sources on $AB$, $BC$, and $CA$.
For stochastic response functions, introduce independent private randomness
at each inflated response; it is suppressed in the following notation.
Take independent identically distributed copies $U_1,U_2,V_1,V_2,W_1,W_2$
and define
\[
\begin{array}{lll}
 A_1=f_A(U_1,W_1),&B_1=f_B(U_1,V_1),&C_1=f_C(V_1,W_1),\\
 A_2=f_A(U_2,W_1),&B_2=f_B(U_1,V_2),&C_2=f_C(V_1,W_2).
\end{array}
\]
Every classical triangle realization therefore induces a joint distribution
$q$ of these six outputs. The following five triple marginals of $q$ are
fixed by the original behavior $p$:
\begin{equation}
\begin{aligned}
 q_{A_1B_1C_1}(a,b,c)&=p_{ABC}(a,b,c),\\
 q_{A_2B_2C_1}(a,b,c)&=p_{AC}(a,c)p_B(b),\\
 q_{A_1B_2C_2}(a,b,c)&=p_{AB}(a,b)p_C(c),\\
 q_{A_2B_1C_2}(a,b,c)&=p_{BC}(b,c)p_A(a),\\
 q_{A_2B_2C_2}(a,b,c)&=p_A(a)p_B(b)p_C(c).
\end{aligned}
\label{eq:spiralmarginals}
\end{equation}
For example, $(A_2,C_1)$ has the original $AC$ marginal, and all its
ancestors are disjoint from those of $B_2$. This proves the second
factorization. The other factorizations follow from the same argument.
The five sets are the maximal ancestrally-independent-expressible sets of
this spiral. Their pair marginals need not be imposed separately.

Fix finite output alphabet sizes $(d_A,d_B,d_C)$ and let
$d=d_Ad_Bd_C$. Order joint assignments lexicographically as
$(A_1,B_1,C_1,A_2,B_2,C_2)$, and order marginal rows by the five blocks in
Eq.~\eqref{eq:spiralmarginals}, then lexicographically within each block.
The induced marginal program is
\begin{equation}
 Aq=b(p),\qquad q\geq0,
 \label{eq:spirallp}
\end{equation}
where $A$ is a $5d\times d^2$ integer matrix containing exactly $5d^2$
nonzero entries, all equal to one. Each column contributes one entry to each
marginal block. Normalization of the first block implies $\sum q=1$.
The implementation guards these three sizes before enumerating
assignments, and stores $A$ sparsely. For binary outputs there are 40 rows,
64 variables, and 320 nonzeros.

\begin{proposition}[Scope of the spiral marginal test]
Every no-input classical triangle behavior makes
Eq.~\eqref{eq:spirallp} feasible. Consequently infeasibility excludes that
resource model. Feasibility of this program alone does not certify triangle
compatibility.
\end{proposition}
\begin{proof}
The copied-source construction produces an admissible nonnegative normalized
$q$, and its marginals satisfy Eq.~\eqref{eq:spiralmarginals}. The converse
would require reconstructing the source model, which the marginal LP does
not enforce.
\end{proof}
The common joint distribution construction relies on classical variables
being reusable at multiple response nodes. It does not justify copying
quantum source systems or assuming joint outcomes for incompatible quantum
operations. This LP therefore provides no general quantum-network
incompatibility claim. Quantum inflation requires a different operator
construction.

\subsection{Exact infeasibility certificates}
A candidate dual vector $y$ is a valid contradiction precisely when
\begin{equation}
 A^{\mathsf T}y\geq0,\qquad y^{\mathsf T}b(p)<0.
 \label{eq:spiralfarkas}
\end{equation}
Indeed, a hypothetical nonnegative $q$ satisfying the marginal equations
would give $0\leq q^{\mathsf T}A^{\mathsf T}y=y^{\mathsf T}b(p)<0$.
Conversely, Farkas' lemma guarantees a separating vector for every infeasible
program. Bounding the coordinates of $y$ during numerical discovery, for
example by $|y_j|\leq1$, loses no strict separation because a Farkas vector
can be scaled. These discovery bounds are not part of certificate validity.

All input probabilities and reconstructed certificate coefficients are
rational. Verification rebuilds $A$ and $b(p)$ and checks every inequality
in Eq.~\eqref{eq:spiralfarkas} exactly, independently of the numerical
solver. A useful exact repair adds $\delta=-\min(A^{\mathsf T}y)>0$ to
every coefficient in the first marginal block whenever a candidate has
negative dual slack. Every slack then increases by $\delta$, and the dual
margin increases by $\delta\sum p_{ABC}=\delta$. The repaired vector is
retained only if its final exact margin is still negative.

A numerically reported infeasibility without a verified dual is not an
exact exclusion. A rational primal $q$ satisfying the equations certifies
only feasibility of the relaxation, so its compatibility status remains
\emph{not ruled out}. Numerical failures, missing or malformed witnesses, and unsuccessful
rational reconstruction retain an explicit unresolved status or numerical
evidence label. Numerical feasibility additionally requires checked residuals;
solver option values that would be silently reset are rejected.

\subsection{A W-type benchmark beyond the Finner checks}
For binary outputs, the event inclusion on the six inflated variables gives
\begin{equation}
 p_A(1)p_B(1)p_C(1)\leq p_{ABC}(000)
 +p_{AC}(11)p_B(1)+p_{AB}(11)p_C(1)+p_{BC}(11)p_A(1).
 \label{eq:spiralw}
\end{equation}
To see the inclusion, if all three outer variables are one, either all
inner variables are zero, or at least one of $A_1,B_1,C_1$ is one. The
latter possibilities lie in the corresponding expressible triple events
on the right. Taking probabilities and applying
Eq.~\eqref{eq:spiralmarginals} proves the inequality.

The distribution $p(100)=p(010)=p(001)=1/3$ has a zero right-hand side and
left-hand side $1/27$. This provides a simple integer dual certificate:
in the stated zero-based row order, set $y_0=y_{15}=y_{23}=y_{31}=1$,
$y_{39}=-1$, and every other coefficient to zero. Direct checking gives
$\min(A^{\mathsf T}y)=0$ and $y^{\mathsf T}b=-1/27$.
Meanwhile every nonzero event passes the triangle Finner test, since
$(1/3)^2\leq(1/3)(2/3)(2/3)$. The benchmark thus exercises
stronger classical compatibility information without changing its evidence
scope. It is an incompatibility target, not a demonstrated quantum advantage:
no quantum realization of this target is supplied by this example.

\subsection{A polynomial witness, not a linear Bell functional}
For fixed output alphabets, $A$ is independent of the tested distribution.
An exactly verified $A^{\mathsf T}y\geq0$ therefore defines the universal
classical-triangle inequality $\Phi_y(p)=y^{\mathsf T}b(p)\geq0$.
Choosing $y$ by optimizing against one distribution does not invalidate
this statement once structural dual feasibility has been independently
checked. However, $b(p)$ contains linear, quadratic, and cubic terms in the
original probabilities. The witness is linear in the lifted marginal
coordinates, but is generally a polynomial in $p$, not an ordinary Bell
score.

Exported numerical-discovery records bind the complete behavior,
source graph, output alphabets, and row-order schema through canonical
fingerprints. Checking the same coefficient vector on different data must
explicitly rebuild and evaluate $b(p)$ in the new context. If a future
relaxation introduces distribution-dependent coefficients in $A$ itself,
the preceding universal-polynomial interpretation requires a new proof;
a single-instance dual cannot then be reused without checking the changed
constraints.
